\documentclass[a4paper,12pt,oneside]{maki-notes}

\renewcommand{\LectureNotesTitle}{Workflow Guide Test}
\renewcommand{\AuthorInfo}{Test Runner}

\SetCourseInfo{
  course={高等数学},
  term={2026 春},
  audience={考研数学},
  instructor={Maki},
  lecture={第 3 讲},
  date={2026-04-13}
}

\begin{document}
\frontmatter
\tableofcontents

\mainmatter

\SetChapterInfo{
  subtitle={从定义到连续性},
  lecture={第 1 讲},
  date={2026-04-13},
  keywords={极限, 连续, 夹逼准则},
  prereq={函数与实数基础},
  goals={理解极限定义; 会判断连续性; 会使用夹逼准则}
}

\chapter{极限与连续}
\MakeChapterGuide{
  route={定义 -> 性质 -> 计算 -> 连续},
  formulas={
    \[
      \lim_{x \to x_0} f(x) = A
    \]
  },
  pitfalls={连续不必可导; 去心邻域条件不能漏写},
  methods={夹逼, 等价无穷小, 代数变形},
  examfocus={极限定义判断题, 连续性与分段函数}
}

\section{函数极限}
这里用来验证章节导读页后的正文和局部导航能否稳定共存。

\begin{keyformula}
\[
  \lim_{x \to x_0} f(x) = A
\]
\end{keyformula}

\begin{pitfall}[易错点]
写极限定义时不能漏掉去心邻域条件。
\end{pitfall}

\section{连续性}
这一节提供第二个 section，确保本章导航不是空结构。

\begin{methodnote}
判断连续性时要同时检查极限存在与函数值匹配。
\end{methodnote}

\begin{examfocus}
分段函数在分界点的连续性与可导性是高频考点。
\end{examfocus}

\chapter{导数}
这一章不设置章节元信息，用来验证上一章的专属字段不会泄漏到下一章。

\section{导数定义}
这里用于验证未设置章节元信息时文档仍能正常编译。

\end{document}
